\documentclass[11pt]{article}

\usepackage{amsmath,amssymb,amsthm,bm,mathtools}
\usepackage{geometry}
\geometry{margin=1in}
\usepackage{setspace}
\usepackage{enumitem}
\usepackage{booktabs}
\usepackage{hyperref}
\usepackage{algorithm}
\usepackage{algpseudocode}
\usepackage{float}

\onehalfspacing

\newcommand{\R}{\mathbb R}
\newcommand{\ind}{\mathbb I}
\newcommand{\norm}[1]{\left\lVert #1 \right\rVert}
\newcommand{\trans}{^{\mathsf T}}
\DeclareMathOperator{\logit}{logit}
\DeclareMathOperator{\IG}{IG}
\DeclareMathOperator{\GammaDist}{Gamma}
\DeclareMathOperator{\Bern}{Bernoulli}
\DeclareMathOperator{\NB}{NB}
\DeclareMathOperator{\EPA}{EPA}
\DeclareMathOperator{\Perm}{Perm}
\DeclareMathOperator{\Unif}{Uniform}

\title{Joint Multilayered LSIRM with the\\
Ewens--Pitman Attraction Pairwise Partition Prior on Items\\
(EPA--LSIRM, v12)}
\date{Model Description}

\begin{document}
\maketitle

\section{Overview}

This document defines a joint multilayered LSIRM in which item clustering is
performed by the Ewens--Pitman Attraction (EPA) pairwise partition prior of
Dahl, Day, and Tsai (2017), conditioned on item latent positions.  The
previous formulation (v11) placed a finite Gaussian mixture (NIW) prior
\emph{generatively} on item latent positions $b_j^{(\ell)}$.  In v12 the
generative mixture is removed: items have a single common Gaussian prior on
their latent positions, and an EPA prior is imposed on the \emph{partition}
of items only, with similarity driven by pairwise distances among those same
positions.

The structural consequence is the following \emph{decoupling property}.  The
LSIRM response likelihood depends on item positions $z_q$ but not on the
partition $\mathcal{P}$, and the prior on $b_j^{(\ell)}$ also does not depend
on $\mathcal{P}$.  Therefore, in any Metropolis--Hastings move that changes
$\mathcal{P}$ alone (single-item Gibbs, split--merge), the LSIRM likelihood
factor $p(Y\mid\Theta)$ cancels exactly in the MH ratio; the acceptance
probability reduces to a ratio of EPA partition pmfs and a proposal
correction.  Conversely, in any move that changes $b_j^{(\ell)}$, the EPA
prior must be re-evaluated because pairwise distances change.

The LSIRM block --- linear predictors, per-layer observation models
(Bernoulli, Student-$t$, Negative Binomial, LSGRM), respondent and item
effects, identifiability conventions --- is identical to v11 and is
reproduced below for self-containedness.  The only LSIRM-side change is in
the prior on $b_j^{(\ell)}$, which is now an isotropic Gaussian rather than a
cluster-specific mixture.  All clustering machinery in §\ref{sec:epa-prior}
through §\ref{sec:sm} is new in v12.

\section{Indexing, dimensions, and observed data}

\begin{center}
\begin{tabular}{lll}
\toprule
Symbol & Meaning & Dimension/range \\
\midrule
$\ell$ & layer index & $\ell=1,\dots,L$ \\
$i$ & respondent index & $i=1,\dots,n_\ell$ \\
$j$ & local item index within layer $\ell$ & $j=1,\dots,P_\ell$ \\
$q$ & global item index across all layers & $q=1,\dots,P$, $P=\sum_{\ell=1}^L P_\ell$ \\
$\ell(q)$ & layer containing global item $q$ & $\{1,\dots,L\}$ \\
$j(q)$ & local item index of global item $q$ & $\{1,\dots,P_{\ell(q)}\}$ \\
$d$ & latent-space dimension & $d=2$ in the current implementation \\
$\mathcal{P}$ & partition of $\{1,\dots,P\}$ & $\{S_1,\dots,S_K\}$ \\
$K$ & number of subsets in $\mathcal{P}$ & $K=|\mathcal{P}|$ \\
$\sigma$ & EPA allocation permutation & $\sigma\in\Perm(P)$ \\
\bottomrule
\end{tabular}
\end{center}

\paragraph{Common-respondent assumption.}
Throughout this document the typical multilayered setting with shared
respondents is assumed: $n_\ell=n$ for all $\ell$, and the respondent latent
position is shared across layers, $a_i^{(\ell)}=a_i\in\R^d$.  The layered
notation $a_i^{(\ell)}$ is kept in equations below for visual symmetry with
$b_j^{(\ell)}$ but should be read as the single $a_i$.  Two consequences
follow.  First, every $b_j^{(\ell)}$ lives in one common coordinate system
anchored by the shared $a$'s; the only identifiability ambiguity is a single
global orthogonal transformation of the entire configuration
$(a, b^{(1)},\dots,b^{(L)})$, not per-layer rigid motions.  Second, the MH
update of $a_i$ aggregates likelihood contributions across all layers
containing respondent $i$, with a single fixed unit-variance prior on $a_i$.

The observed data are
\[
    Y^{(\ell)}=\{y_{ij}^{(\ell)}:(i,j)\in\mathcal O_\ell\},
\]
with $\mathcal O_\ell$ the set of observed respondent--item pairs in layer
$\ell$.

\section{LSIRM likelihood}

The model uses, for respondent $i$ and item $j$ in layer $\ell$:
\begin{itemize}[leftmargin=1.5em]
    \item $a_i\in\R^d$: shared respondent latent position.
    \item $b_j^{(\ell)}\in\R^d$: item latent position in layer $\ell$.
    \item $\alpha_i^{(\ell)}\in\R$: respondent effect in layer $\ell$.
    \item $\beta_j^{(\ell)}\in\R$: item difficulty for binary, continuous,
    and count layers; replaced by a threshold vector
    $\beta_{j,1}^{(\ell)}<\cdots<\beta_{j,K_\ell-1}^{(\ell)}$ in ordinal
    layers.
    \item $\gamma_\ell>0$: layer-specific distance weight.
    \item $\sigma_0^2>0$: continuous-layer error scale.
    \item $\nu_2>0$: continuous-layer Student-$t$ degrees of freedom.
    \item $\kappa_j>0$: item-specific count-layer overdispersion.
\end{itemize}
Let $\mathcal L_{\mathrm{ord}}\subseteq\{1,\dots,L\}$ index the ordinal
layers.

\paragraph{Implementation note.}
The current implementation is a four-layer instance ($L=4$): one binary, one
continuous, one count, and one ordinal (LSGRM) layer, so
$\mathcal L_{\mathrm{ord}}=\{4\}$.  The notation below is layer-generic.

The layer-specific LSIRM linear predictor is
\begin{equation}
    \eta_{ij}^{(\ell)}
    =
    \alpha_i^{(\ell)}-\beta_j^{(\ell)}-\gamma_\ell\norm{a_i-b_j^{(\ell)}}_2,
    \qquad \ell\notin\mathcal L_{\mathrm{ord}},
    \label{eq:lsirm-linear-predictor}
\end{equation}
with the convention that for ordinal layers the difficulty $\beta_j^{(\ell)}$
is replaced by item-specific thresholds entering the likelihood directly:
\begin{equation}
    \eta_{ij}^{(\ell)}
    =
    \alpha_i^{(\ell)}-\gamma_\ell\norm{a_i-b_j^{(\ell)}}_2,
    \qquad \ell\in\mathcal L_{\mathrm{ord}}.
    \label{eq:lsirm-linear-predictor-ord}
\end{equation}

\subsection{Per-layer observation models}

\paragraph{Binary (logistic regression).}
\begin{equation}
    y_{ij}^{(\ell)}\mid \eta_{ij}^{(\ell)}\sim
    \Bern(p_{ij}^{(\ell)}),
    \qquad \logit p_{ij}^{(\ell)}=\eta_{ij}^{(\ell)}.
\end{equation}

\paragraph{Continuous (robust Student-$t$ via normal--gamma scale mixture).}
\begin{align}
    y_{ij}^{(\ell)}\mid \eta_{ij}^{(\ell)},\lambda_{ij}^{(\ell)}
    &\sim N\!\left(\eta_{ij}^{(\ell)},\;\sigma_0^2/\lambda_{ij}^{(\ell)}\right),\notag\\
    \lambda_{ij}^{(\ell)}
    &\sim \GammaDist(\nu_2/2,\nu_2/2).
\end{align}
Marginalizing $\lambda_{ij}^{(\ell)}$ gives
$y_{ij}^{(\ell)}\sim t_{\nu_2}(\eta_{ij}^{(\ell)},\sigma_0^2)$.

\paragraph{Count (negative binomial with item-specific overdispersion).}
\begin{equation}
    y_{ij}^{(\ell)}\mid \eta_{ij}^{(\ell)},\kappa_j
    \sim \NB\!\left(r=1/\kappa_j,\;\mu=\exp\eta_{ij}^{(\ell)}\right),
\end{equation}
with $E[y_{ij}^{(\ell)}]=\mu$ and $\mathrm{Var}[y_{ij}^{(\ell)}]=\mu+\kappa_j\mu^2$.
The Poisson limit is recovered as $\kappa_j\to 0$.  The prior is log-normal,
\begin{equation}
    \log\kappa_j\sim N(\mu_{\log\kappa},\sigma_{\log\kappa}^2).
\end{equation}

\paragraph{Ordinal (LSGRM, De Carolis--Kang--Jeon, 2025).}
For an ordinal layer $\ell\in\mathcal L_{\mathrm{ord}}$ with $K_\ell$ ordered
categories and item thresholds
$\beta_{j,1}^{(\ell)}<\cdots<\beta_{j,K_\ell-1}^{(\ell)}$,
\begin{equation}
    \logit P\!\bigl(y_{ij}^{(\ell)}\geq k+1\bigm|\eta_{ij}^{(\ell)},\beta_{j,k}^{(\ell)}\bigr)
    =\eta_{ij}^{(\ell)}-\beta_{j,k}^{(\ell)},
    \qquad k=1,\dots,K_\ell-1,
\end{equation}
with $P(y\geq 1)=1$, $P(y\geq K_\ell+1)=0$; category probabilities follow by
differencing.

\subsection{Compact notation}

Each per-layer likelihood is summarized as
\begin{equation}
    y_{ij}^{(\ell)}\mid \eta_{ij}^{(\ell)},\xi_\ell
    \sim f_\ell\bigl\{y_{ij}^{(\ell)}\mid \eta_{ij}^{(\ell)},\xi_\ell\bigr\},
    \qquad (i,j)\in\mathcal O_\ell,
\end{equation}
where $\xi_\ell$ collects layer-specific nuisance parameters (e.g.\
$(\sigma_0^2,\lambda_{ij}^{(\ell)})$ in continuous layers,
$\{\kappa_j\}$ in count layers, the threshold vector in ordinal layers, and
$\gamma_\ell$ in every layer).  The full LSIRM log-likelihood is
\begin{equation}
    \ell_{\mathrm{LSIRM}}
    =
    \sum_{\ell=1}^L\sum_{(i,j)\in\mathcal O_\ell}
    \log f_\ell\bigl\{y_{ij}^{(\ell)}\mid \eta_{ij}^{(\ell)},\xi_\ell\bigr\}.
    \label{eq:lsirm-loglik}
\end{equation}

\section{EPA pairwise partition prior on items}
\label{sec:epa-prior}

\subsection{Global item latent positions}

The global item latent position is
\begin{equation}
    z_q = b_{j(q)}^{(\ell(q))}\in\R^d,
    \qquad q=1,\dots,P,
\end{equation}
and we write $z=(z_1,\dots,z_P)$.  The clustering model is placed on the
partition of $\{1,\dots,P\}$, conditional on $z$.

\subsection{Pairwise similarity}

For an EPA temperature $\tau\geq 0$,
\begin{equation}
    \lambda_{qr}(z,\tau) = \exp\!\bigl\{-\tau\,d_{qr}(z)\bigr\},
    \qquad d_{qr}(z)=\norm{z_q-z_r}_2^2.
    \label{eq:lambda}
\end{equation}
Unlike the original Dahl--Day--Tsai (2017) formulation, no latent-scale
normalizer $s_z^2$ is used.  In EPA--LSIRM the latent space scale is already
anchored by the priors $a_i\sim N_d(0,I_d)$ (\ref{eq:a-prior}) and
$b_j^{(\ell)}\sim N_d(0,\sigma_b^2 I_d)$ (\ref{eq:b-prior}), so $\norm{z_q-z_r}_2^2$
is automatically $O(1)$ across layers.  An additional normalizer would only
introduce a redundant reparameterization
$\tau\mapsto\tau/s_z^2$ without altering the model.

\subsection{Sequential allocation EPA pmf}

Let $\sigma=(\sigma_1,\dots,\sigma_P)\in\Perm(P)$ be an allocation permutation
and $\pi_t=\pi(\sigma_1,\dots,\sigma_t)$ the partition produced after the
first $t$ allocations, with $q_{t-1}=|\pi_{t-1}|$.  The conditional
allocation probability of $\sigma_t$ at step $t$ is
\begin{equation}
\Pr\!\bigl(\sigma_t\in S\,\big|\,\sigma_{1:t-1},\alpha,\delta,\tau,z\bigr)
=
\begin{cases}
\displaystyle
\frac{t-1-\delta\,q_{t-1}}{\alpha+t-1}\cdot
\frac{\sum_{r\in S}\lambda_{\sigma_t r}(z,\tau)}
     {\sum_{s=1}^{t-1}\lambda_{\sigma_t \sigma_s}(z,\tau)},
& S\in\pi_{t-1},\\[10pt]
\displaystyle\frac{\alpha+\delta\,q_{t-1}}{\alpha+t-1},
& S \text{ a new subset.}
\end{cases}
\label{eq:epa-cond}
\end{equation}
The full EPA pmf factorizes as
\begin{equation}
    p_{\EPA}(\mathcal{P}\mid\sigma,z,\alpha,\delta,\tau)
    =\prod_{t=1}^{P}p_t(\alpha,\delta,\sigma,z,\pi_{t-1}).
    \label{eq:epa-pmf}
\end{equation}
Mass parameter $\alpha\in\R_{>0}$, discount $\delta\in[0,1)$, and temperature
$\tau\geq 0$.

\subsection{Hyperpriors and identifiability conventions}

\begin{align}
    \alpha &\sim \GammaDist(a_\alpha,b_\alpha),\\
    \tau &\sim \GammaDist(a_\tau,b_\tau),\\
    \sigma &\sim \Unif(\Perm(P)),\\
    \delta &= 0\quad\text{(fixed in the current implementation).}
\end{align}
Defaults are weakly informative: $(a_\alpha,b_\alpha)=(1,1)$,
$(a_\tau,b_\tau)=(1,1)$.  Setting $\delta=0$ reduces the EPA mass-weighting
factor $(\alpha+\delta q_{t-1})/(\alpha+t-1)$ to $\alpha/(\alpha+t-1)$,
matching the Dirichlet-process--type sequential allocation under EPA's
distance-dependent reallocation kernel.

\subsection{Decoupling property}
\label{sec:decouple}

The partition $\mathcal{P}$ enters the joint distribution \emph{only} through
the EPA pmf \eqref{eq:epa-pmf}.  The LSIRM likelihood
\eqref{eq:lsirm-loglik} does not depend on $\mathcal{P}$, and no other prior
in the model depends on $\mathcal{P}$.  Therefore
\begin{equation}
    p\bigl(Y,\mathcal{P}\,\big|\,\Theta_{\mathrm{LSIRM}},\sigma,\alpha,\delta,\tau\bigr)
    =
    p(Y\mid\Theta_{\mathrm{LSIRM}})\;
    p_{\EPA}(\mathcal{P}\mid\sigma,z,\alpha,\delta,\tau),
    \label{eq:joint-decouple}
\end{equation}
which is the structural property that drives all partition updates below.

\section{Priors for non-clustering parameters}

\paragraph{Respondent latent position (fixed unit variance for identifiability).}
\begin{equation}
    a_i\sim N_d(0,I_d),\qquad i=1,\dots,n.
    \label{eq:a-prior}
\end{equation}
Fixing the variance at $1$ anchors the latent-space scale; following LSGRM
convention, $\gamma_\ell$ remains the free distance-scale parameter.

\paragraph{Item latent position (free isotropic Gaussian prior, modified from v11).}
In contrast to v11, where $b_j^{(\ell)}$ inherited its prior from a
cluster-specific Gaussian mixture, v12 endows each $b_j^{(\ell)}$ with an
independent isotropic Gaussian prior of fixed scale,
\begin{equation}
    b_j^{(\ell)}\sim N_d(0,\sigma_b^2 I_d),
    \qquad j=1,\dots,P_\ell,\;\ell=1,\dots,L.
    \label{eq:b-prior}
\end{equation}
The scale $\sigma_b^2$ is a fixed hyperparameter (not sampled); the
implementation uses $\sigma_b^2=1$ to match the respondent prior scale and to
keep $\gamma_\ell$ as the only free distance-scale parameter.  The cluster
identity of item $j$ does \emph{not} enter \eqref{eq:b-prior}, by design ---
this is what produces the decoupling property in §\ref{sec:decouple}.

\paragraph{Respondent effect.}
\begin{equation}
    \alpha_i^{(\ell)}\mid\sigma_{\alpha,\ell}^2\sim N(0,\sigma_{\alpha,\ell}^2),
    \qquad
    \sigma_{\alpha,\ell}^2\sim\IG(A_{\alpha,\ell},B_{\alpha,\ell}),
    \qquad \ell=1,\dots,L.
    \label{eq:alpha-prior}
\end{equation}

\paragraph{Item difficulty $\beta_j^{(\ell)}$ (fixed-scale prior).}
For non-ordinal layers,
\begin{equation}
    \beta_j^{(\ell)}\sim N(0,\sigma_{\beta,\ell}^2),
    \qquad \ell\notin\mathcal L_{\mathrm{ord}},
    \label{eq:beta-prior}
\end{equation}
with $\sigma_{\beta,\ell}^2$ fixed (the implementation uses
$\sigma_{\beta,\ell}=2$).

\paragraph{Ordinal thresholds (independent normals subject to ordering).}
For $\ell\in\mathcal L_{\mathrm{ord}}$,
\begin{equation}
    \beta_{j,k}^{(\ell)}\sim N(0,\sigma_{\beta,\ell}^2),
    \qquad
    \beta_{j,1}^{(\ell)}<\cdots<\beta_{j,K_\ell-1}^{(\ell)},
    \label{eq:beta-thr-prior}
\end{equation}
with $\sigma_{\beta,\ell}^2$ fixed.

The remaining layer-specific nuisance parameters --- $\gamma_\ell$,
$\sigma_0^2$, and $\{\kappa_j\}$ --- are equipped with independent log-normal
or inverse-gamma priors as in v11; their priors are unchanged.

\section{Rotation and reflection invariance}
\label{sec:invariance}

Under the common-respondent assumption, the entire configuration
$(a,b^{(1)},\dots,b^{(L)})$ shares one coordinate system.  The only
identifiability ambiguity is a single global orthogonal transformation
$Q\in O(d)$:
\begin{equation}
    a_i\mapsto Qa_i,
    \qquad
    b_j^{(\ell)}\mapsto Qb_j^{(\ell)}.
\end{equation}
The LSIRM linear predictor depends on positions only through
$\norm{a_i-b_j^{(\ell)}}_2$, which is $O(d)$-invariant.  The EPA similarity
\eqref{eq:lambda} depends only on pairwise distances $\norm{z_q-z_r}_2$ and
is therefore also $O(d)$-invariant.  The priors $N_d(0,I_d)$ on $a_i$ and
$N_d(0,\sigma_b^2 I_d)$ on $b_j^{(\ell)}$ are $O(d)$-invariant in
distribution.  Therefore the joint posterior on the clustering targets
($\mathcal{P}$, $K=|\mathcal{P}|$, the cluster sizes, the posterior
similarity matrix $C_{qr}=\Pr(q\sim_{\mathcal{P}}r\mid Y)$) is $O(d)$-invariant.

\paragraph{Consequence for the sampler.}
Within-MCMC orientation alignment is \emph{not} required for clustering
correctness.  A post-hoc Procrustes alignment is recommended only when the
visualization of $a$ or $b^{(\ell)}$ must be compared to a target
configuration.  Translation invariance is broken by the centered Gaussian
priors on $a$ and $b$, exactly as in v11.

\section{Full joint posterior}

Let
\begin{align*}
    \Theta_{\mathrm{LSIRM}}
    &=\{a_i,\alpha_i^{(\ell)},b_j^{(\ell)},\beta_j^{(\ell)},\xi_\ell,\sigma_{\alpha,\ell}^2:i,j,\ell\},\\
    \Theta_{\mathrm{EPA}}
    &=\{\mathcal{P},\sigma,\alpha,\tau\}\quad(\delta\equiv 0).
\end{align*}
Then
\begin{align}
&p\bigl(\Theta_{\mathrm{LSIRM}},\Theta_{\mathrm{EPA}}\,\big|\,Y\bigr)
\notag\\
&\quad\propto
\prod_{\ell=1}^L
\prod_{(i,j)\in\mathcal{O}_\ell}
    f_\ell\bigl\{y_{ij}^{(\ell)}\,\big|\,\eta_{ij}^{(\ell)},\xi_\ell\bigr\}
\notag\\
&\qquad\times\;
p_{\EPA}(\mathcal{P}\mid\sigma,z,\alpha,\delta,\tau)
\notag\\
&\qquad\times\;
\prod_{i=1}^n N_d(a_i;0,I_d)\;
\prod_{\ell=1}^L\prod_{j=1}^{P_\ell} N_d\bigl(b_j^{(\ell)};0,\sigma_b^2 I_d\bigr)
\notag\\
&\qquad\times\;
\prod_{\ell=1}^L\!\left[
\prod_{i=1}^n p\bigl(\alpha_i^{(\ell)}\mid\sigma_{\alpha,\ell}^2\bigr)\,
p_\beta^{(\ell)}\bigl(\beta^{(\ell)}\bigr)\,
p(\sigma_{\alpha,\ell}^2)\,p(\xi_\ell)
\right]
\notag\\
&\qquad\times\;
p(\alpha)\,p(\tau)\,p(\sigma),
\label{eq:full-posterior}
\end{align}
with $z_q=b_{j(q)}^{(\ell(q))}$ and the layer-specific item-difficulty prior
factor $p_\beta^{(\ell)}(\beta^{(\ell)})$ defined as in
\eqref{eq:beta-prior}--\eqref{eq:beta-thr-prior}.

\section{Metropolis--Hastings updates for LSIRM parameters}

All LSIRM parameters with non-conjugate full conditionals are updated by
random-walk Metropolis--Hastings with symmetric Gaussian proposals.

\subsection{Respondent effect $\alpha_i^{(\ell)}$}

Propose $\alpha_i^{(\ell)\prime}=\alpha_i^{(\ell)}+\varepsilon_\alpha$,
$\varepsilon_\alpha\sim N(0,s_\alpha^2)$.  The log MH ratio is
\begin{align}
    \log R_\alpha
    &=
    \sum_{j:(i,j)\in\mathcal{O}_\ell}
    \!\left[
    \log f_\ell\bigl\{\cdot\mid\eta_{ij}^{(\ell)\prime}\bigr\}-
    \log f_\ell\bigl\{\cdot\mid\eta_{ij}^{(\ell)}\bigr\}
    \right]
    \notag\\
    &\quad-\frac{1}{2\sigma_{\alpha,\ell}^2}
    \bigl\{(\alpha_i^{(\ell)\prime})^2-(\alpha_i^{(\ell)})^2\bigr\}.
    \label{eq:alpha-mh-ratio}
\end{align}

\subsection{Item difficulty $\beta_j^{(\ell)}$}

For non-ordinal layers,
\begin{align}
    \log R_\beta
    &=
    \sum_{i:(i,j)\in\mathcal{O}_\ell}
    \!\left[
    \log f_\ell\bigl\{\cdot\mid\eta_{ij}^{(\ell)\prime}\bigr\}-
    \log f_\ell\bigl\{\cdot\mid\eta_{ij}^{(\ell)}\bigr\}
    \right]
    \notag\\
    &\quad
    -\frac{1}{2\sigma_{\beta,\ell}^2}\bigl\{(\beta_j^{(\ell)\prime})^2-(\beta_j^{(\ell)})^2\bigr\}.
    \label{eq:beta-mh-ratio}
\end{align}
For ordinal layers, propose each threshold by an independent Gaussian random
walk, reject any move that violates ordering, and use the analogue of
\eqref{eq:beta-mh-ratio} with the LSGRM ordinal likelihood evaluated at the
full threshold vector.

\subsection{Respondent latent position $a_i$}

Propose $a_i'=a_i+\varepsilon_a$, $\varepsilon_a\sim N_d(0,s_a^2 I_d)$, and
aggregate likelihood contributions from every layer:
\begin{align}
    \log R_a
    &=
    \sum_{\ell=1}^L
    \sum_{j:(i,j)\in\mathcal{O}_\ell}
    \!\left[
    \log f_\ell\bigl\{\cdot\mid\eta_{ij}^{(\ell)\prime}\bigr\}-
    \log f_\ell\bigl\{\cdot\mid\eta_{ij}^{(\ell)}\bigr\}
    \right]
    \notag\\
    &\quad-\frac{1}{2}\bigl\{\norm{a_i'}_2^2-\norm{a_i}_2^2\bigr\}.
    \label{eq:a-mh-ratio}
\end{align}

\subsection{Item latent position $b_j^{(\ell)}$ (modified from v11)}
\label{sec:b-update}

Let $q=q(\ell,j)$.  Propose $b_j^{(\ell)\prime}=b_j^{(\ell)}+\varepsilon_b$,
$\varepsilon_b\sim N_d(0,s_b^2 I_d)$, and let $z'$ denote $z$ with $z_q$
replaced by $b_j^{(\ell)\prime}$.  The log MH ratio combines three terms ---
LSIRM likelihood, free Gaussian prior on $b$, \emph{and} the EPA prior, the
last of which is required because the EPA pmf depends on $z$ through pairwise
similarity:
\begin{align}
    \log R_b
    &=
    \sum_{i:(i,j)\in\mathcal{O}_\ell}
    \!\left[
    \log f_\ell\bigl\{\cdot\mid\eta_{ij}^{(\ell)\prime}\bigr\}-
    \log f_\ell\bigl\{\cdot\mid\eta_{ij}^{(\ell)}\bigr\}
    \right]
    \notag\\
    &\quad
    -\frac{1}{2\sigma_b^2}\bigl\{\norm{b_j^{(\ell)\prime}}_2^2-\norm{b_j^{(\ell)}}_2^2\bigr\}
    \notag\\
    &\quad
    +\log p_{\EPA}(\mathcal{P}\mid\sigma,z',\alpha,\delta,\tau)
    -\log p_{\EPA}(\mathcal{P}\mid\sigma,z,\alpha,\delta,\tau).
    \label{eq:b-mh-ratio}
\end{align}
This replaces the v11 mixture term: the cluster identity $c_q$ no longer
appears, but the EPA prior contribution does, and is non-trivial because
moving $z_q$ rescales similarities $\lambda_{qr}(z,\tau)$ for every $r\ne q$.

\paragraph{Caching.}
Recomputing \eqref{eq:epa-pmf} from scratch costs $O(P^2)$.  Updating one
$z_q$ changes only the $q$-th row and column of the similarity matrix
$\Lambda$; cached partial sums $\sum_{r\in S}\lambda_{\sigma_t r}$ for each
allocation step $t$ that involves $q$ should be updated incrementally,
giving $O(P)$ work per $b$ update in the typical case.

\subsection{Conjugate update for $\sigma_{\alpha,\ell}^2$}

\begin{equation}
    \sigma_{\alpha,\ell}^2\mid\alpha^{(\ell)}\sim
    \IG\!\left(A_{\alpha,\ell}+\tfrac{n}{2},\;
    B_{\alpha,\ell}+\tfrac{1}{2}\textstyle\sum_i(\alpha_i^{(\ell)})^2\right).
    \label{eq:sigma-alpha-update}
\end{equation}

\section{Single-item Gibbs update for $\mathcal{P}$}

For each global item $q\in\{1,\dots,P\}$ in random order, remove $q$ from its
current subset and reassign it to either an existing subset
$S\in\mathcal{P}_{-q}$ or a new singleton.  Let $\mathcal{P}^{(S)}$ denote
the candidate partition obtained by adding $q$ to $S$ (or creating a new
singleton).  Because of the decoupling \eqref{eq:joint-decouple} the LSIRM
likelihood does not enter, and the conditional is
\begin{equation}
    \Pr(q\in S\mid\mathcal{P}_{-q},\sigma,z,\alpha,\delta,\tau)
    \propto p_{\EPA}\bigl(\mathcal{P}^{(S)}\mid\sigma,z,\alpha,\delta,\tau\bigr).
    \label{eq:single-item-gibbs}
\end{equation}
The candidates are: each existing subset of $\mathcal{P}_{-q}$, plus a new
singleton, giving $|\mathcal{P}_{-q}|+1$ options.  Because the EPA pmf is
non-exchangeable, each candidate requires re-evaluating
\eqref{eq:epa-pmf} for the corresponding partition --- but since adding a
single item only changes the EPA factor at the position $t$ where $q$ appears
in $\sigma$ and at later steps, the per-candidate cost is reduced by caching
the cumulative log-products of $p_t$ for $t$ preceding the modified position.

\section{Permutation update}

Update $\sigma$ by random-swap MH: propose $\sigma'$ obtained from $\sigma$
by swapping two uniformly chosen positions, accept with
\begin{equation}
    \log R_\sigma
    =\log p_{\EPA}(\mathcal{P}\mid\sigma',z,\alpha,\delta,\tau)
    -\log p_{\EPA}(\mathcal{P}\mid\sigma,z,\alpha,\delta,\tau).
    \label{eq:sigma-mh-ratio}
\end{equation}
The uniform prior on $\Perm(P)$ contributes equal terms that cancel.  Several
swaps may be performed per sweep ($M_\sigma$ swaps; default $M_\sigma=P$).

\section{EPA hyperparameter updates}

Update $\alpha$ on the log scale: propose $\log\alpha'=\log\alpha+\varepsilon$,
$\varepsilon\sim N(0,s_{\log\alpha}^2)$, with log-acceptance
\begin{align}
    \log R_\alpha^{\mathrm{hyp}}
    &=\log p_{\EPA}(\mathcal{P}\mid\sigma,z,\alpha',\delta,\tau)
    -\log p_{\EPA}(\mathcal{P}\mid\sigma,z,\alpha,\delta,\tau)
    \notag\\
    &\quad+\log p_\alpha(\alpha')-\log p_\alpha(\alpha)
    +\log\alpha'-\log\alpha,
    \label{eq:alpha-hyp-mh-ratio}
\end{align}
the last line being the log-Jacobian of the log-scale parameterization.  The
update for $\tau$ is analogous.

\section{Jain--Neal restricted Gibbs split--merge}
\label{sec:sm}

The single-item Gibbs sweep on $\mathcal{P}$ in
\eqref{eq:single-item-gibbs} mixes locally but cannot relocate an entire
block of items in one move.  When two distinct clusterings of nearly equal
posterior mass are separated by a barrier of single-item moves through
low-probability intermediate partitions, the chain mixes slowly.  We
therefore add a Jain--Neal split--merge step in the nonconjugate variant of
Jain and Neal (2007), specialized to EPA via the decoupling
\eqref{eq:joint-decouple}.

\subsection{Anchor selection and move type}

Sample distinct anchors $(u,v)$ uniformly without replacement from
$\{1,\dots,P\}^2$ (probability $1/[P(P-1)]$).  Determine move type:
\begin{itemize}
\item If there exists $S\in\mathcal{P}$ with $u,v\in S$, propose a
  \emph{split} of $S$.
\item Otherwise $u\in S_u\ne S_v\ni v$, and propose a \emph{merge} of $S_u$
  and $S_v$.
\end{itemize}
Move type is a deterministic function of $(u,v,\mathcal{P})$, so the
move-type indicator cancels.  Define
\begin{equation}
    \mathcal{S}=\begin{cases}S,&\text{split,}\\ S_u\cup S_v,&\text{merge,}\end{cases}
    \qquad\mathcal{S}^*=\mathcal{S}\setminus\{u,v\}.
\end{equation}

\subsection{Restricted Gibbs sweeps}

For a restricted partition
$\mathbf{L}=\{S^{\mathrm{spl}}_u,S^{\mathrm{spl}}_v\}$ on $\mathcal{S}$ (with
$u\in S^{\mathrm{spl}}_u$, $v\in S^{\mathrm{spl}}_v$ permanently) and a
non-anchor item $q^*\in\mathcal{S}^*$, the conditional reassignment
probability is
\begin{equation}
    \Pr\bigl(L_{q^*}=u\,\big|\,\mathbf{L}_{-q^*},\sigma,z,\alpha,\delta,\tau\bigr)
    =\frac{\rho_u(q^*)}{\rho_u(q^*)+\rho_v(q^*)},
    \label{eq:rgibbs-cond}
\end{equation}
where, after temporarily setting $L_{q^*}=w\in\{u,v\}$ to obtain
$\mathbf{L}^{(w)}$ and combining with the unchanged outside-$\mathcal{S}$
labels,
\begin{equation}
    \rho_w(q^*)
    =p_{\EPA}\bigl(\mathbf{L}^{(w)}\cup\mathcal{P}_{-\mathcal{S}}\,\big|\,\sigma,z,\alpha,\delta,\tau\bigr).
    \label{eq:rho-defn}
\end{equation}
The LSIRM likelihood does not enter \eqref{eq:rho-defn} because of
\eqref{eq:joint-decouple}.

\paragraph{Launch state.}
Initialize a restricted partition on $\mathcal{S}$:
\begin{itemize}
\item \textbf{Split case.} Each $q^*\in\mathcal{S}^*$ is independently
  assigned to $S^{\mathrm{spl}}_u$ or $S^{\mathrm{spl}}_v$ with probability
  $1/2$.
\item \textbf{Merge case.} The trivial restricted partition
  $\{S_u,S_v\}$ is used (the partition we want to merge from, also serving as
  the ``before-launch'' state for the reverse split proposal).
\end{itemize}
Run $R$ restricted Gibbs scans (each scan visits every $q^*\in\mathcal{S}^*$
in fixed order and updates $L_{q^*}$ from \eqref{eq:rgibbs-cond}); denote
the resulting state $\mathbf{L}=\mathbf{L}_R$.  This launch state is
\emph{not} a sample from any target distribution; it serves only as the
starting point of a final scan whose path probability is the proposal
density.

\subsection{Forward proposal (split case)}

From $\mathbf{L}$, run one final restricted Gibbs scan and record the path
probability
\begin{equation}
    q_{\mathrm{prop}}\bigl(\mathcal{P}\to\mathcal{P}^{\mathrm{spl}}\bigr)
    =\prod_{q^*\in\mathcal{S}^*}
        \Pr\bigl(L_{q^*}^{\mathrm{final}}
                =L_{q^*}^{\mathcal{P}^{\mathrm{spl}}}
                \,\big|\,\text{scan-state at $q^*$}\bigr).
    \label{eq:qfwd-split}
\end{equation}
The result defines
\begin{equation}
    \mathcal{P}^{\mathrm{spl}}
    =\bigl(\mathcal{P}\setminus\{S\}\bigr)\cup\{S^{\mathrm{spl}}_u,S^{\mathrm{spl}}_v\}.
\end{equation}
The reverse move is the deterministic merge, so
$q_{\mathrm{prop}}(\mathcal{P}^{\mathrm{spl}}\to\mathcal{P})=1$.

\subsection{Forward proposal (merge case)}

The forward move is deterministic:
\begin{equation}
    \mathcal{P}^{\mathrm{mrg}}
    =\bigl(\mathcal{P}\setminus\{S_u,S_v\}\bigr)\cup\{S_u\cup S_v\},
    \qquad q_{\mathrm{prop}}\bigl(\mathcal{P}\to\mathcal{P}^{\mathrm{mrg}}\bigr)=1.
\end{equation}
The reverse split must recreate $\{S_u,S_v\}$ on $\mathcal{S}=S_u\cup S_v$:
generate a reverse launch state by $R$ restricted Gibbs scans starting from
$\mathbf{L}_0^{\mathrm{rev}}=\{S_u,S_v\}$, then \emph{score} (do not sample)
the original assignment in $\mathcal{P}$ via one scoring scan,
\begin{equation}
    q_{\mathrm{prop}}\bigl(\mathcal{P}^{\mathrm{mrg}}\to\mathcal{P}\bigr)
    =\prod_{q^*\in\mathcal{S}^*}
        \Pr\bigl(L_{q^*}^{\mathrm{final}}
                =L_{q^*}^{\mathcal{P}}
                \,\big|\,\text{reverse scan-state at $q^*$}\bigr).
    \label{eq:qrev-merge}
\end{equation}

\subsection{Acceptance probabilities}

Because of \eqref{eq:joint-decouple} the LSIRM likelihood cancels, the
move-type indicator is deterministic, and the symmetric anchor probability
$1/[P(P-1)]$ also cancels.  Therefore
\begin{align}
    \log R_{\mathrm{spl}}
    &=\log p_{\EPA}(\mathcal{P}^{\mathrm{spl}}\mid\sigma,z,\alpha,\delta,\tau)
    -\log p_{\EPA}(\mathcal{P}\mid\sigma,z,\alpha,\delta,\tau)
    \notag\\
    &\quad
    -\log q_{\mathrm{prop}}(\mathcal{P}\to\mathcal{P}^{\mathrm{spl}}),
    \label{eq:logR-split}\\
    \log R_{\mathrm{mrg}}
    &=\log p_{\EPA}(\mathcal{P}^{\mathrm{mrg}}\mid\sigma,z,\alpha,\delta,\tau)
    -\log p_{\EPA}(\mathcal{P}\mid\sigma,z,\alpha,\delta,\tau)
    \notag\\
    &\quad
    +\log q_{\mathrm{prop}}(\mathcal{P}^{\mathrm{mrg}}\to\mathcal{P}),
    \label{eq:logR-merge}
\end{align}
with acceptance probabilities $\min\{1,\exp(\log R_{\cdot})\}$.

\paragraph{Permutation handling.}
$\sigma$ is held fixed during a split--merge attempt and updated separately
in the random-swap step \eqref{eq:sigma-mh-ratio}.  Because $\sigma$ enters
$p_{\EPA}$ and the restricted-Gibbs conditionals \eqref{eq:rho-defn}
deterministically and consistently across forward and reverse moves, this is
sufficient for detailed balance.

\paragraph{Tuning.}
$R$ controls launch-state quality: $R\in\{2,5,10\}$ following Jain and Neal
(2007); the implementation default is $R=5$.  Combined split/merge
acceptance is targeted at $0.10$--$0.30$.  Per outer sweep, attempt
$M_{\mathrm{SM}}$ split--merge moves (default $M_{\mathrm{SM}}\in\{1,5\}$).

\section{Posterior cluster summaries}

Because cluster labels are not identifiable (subsets in $\mathcal{P}$ are
unordered), summaries should be label-switching invariant.  The primary
summary is the posterior similarity matrix
\begin{equation}
    C_{qr}=\Pr\bigl(q\sim_{\mathcal{P}}r\,\big|\,Y\bigr),
    \qquad q,r=1,\dots,P,
\end{equation}
updated online during MCMC by
$C_{qr}^{\mathrm{raw}}\leftarrow C_{qr}^{\mathrm{raw}}+\ind(q\sim_{\mathcal{P}^{(t)}}r)$
after burn-in and divided by the number of saved iterations at the end.
Also monitor the trace of $K^{(t)}=|\mathcal{P}^{(t)}|$.

\section{One full MCMC sweep}

\begin{algorithm}[H]
\caption{One restricted Gibbs scan
$\textsc{RGibbsScan}(\mathbf{L},\mathcal{S}^*,u,v,\sigma,z,\alpha,\delta,\tau,\mathcal{P}_{-\mathcal{S}})$.
Returns the updated restricted partition and the log-probability of the
sampled path (or, in scoring mode, of a forced path).}
\small
\begin{algorithmic}[1]
\State $\log q\gets 0$
\For{$q^*\in\mathcal{S}^*$ in fixed order}
    \State Compute $\rho_u(q^*),\rho_v(q^*)$ via \eqref{eq:rho-defn}.
    \State $p^*\gets\rho_u(q^*)/(\rho_u(q^*)+\rho_v(q^*))$.
    \If{sampling mode}
        \State Draw $w\in\{u,v\}$ with $\Pr(w=u)=p^*$.
    \Else{ (scoring mode)}
        \State Set $w$ to the prescribed target label $L_{q^*}^{\mathrm{target}}$.
    \EndIf
    \State $L_{q^*}\gets w$;
    update $\mathbf{L}$ accordingly.
    \State $\log q\gets\log q+\log\Pr(L_{q^*}=w)$.
\EndFor
\State \Return $(\mathbf{L},\log q)$.
\end{algorithmic}
\end{algorithm}

\begin{algorithm}[H]
\caption{One split--merge MH step for EPA--LSIRM.}
\small
\begin{algorithmic}[1]
\Require Current state $(\mathcal{P},\sigma,\alpha,\tau,z)$ and number of
launch scans $R\geq 0$.
\State Sample anchors $(u,v)$ uniformly without replacement.
\If{$u\sim_{\mathcal{P}}v$}
    \State $\mathcal{S}\gets S$ (the cluster containing $u,v$);
    \textbf{type}$\gets$\textsc{Split}.
\Else
    \State $\mathcal{S}\gets S_u\cup S_v$; \textbf{type}$\gets$\textsc{Merge}.
\EndIf
\State $\mathcal{S}^*\gets\mathcal{S}\setminus\{u,v\}$.
\If{\textbf{type}$=$\textsc{Split}}
    \State Initialize $\mathbf{L}_0$ by random $1/2$ assignment.
    \For{$r=1,\dots,R$}
        \State $(\mathbf{L}_r,\cdot)\gets\textsc{RGibbsScan}(\mathbf{L}_{r-1},\dots)$ in sampling mode.
    \EndFor
    \State $(\mathbf{L}^{\mathrm{fin}},\log q_{\mathrm{prop}}(\mathcal{P}\to\mathcal{P}^{\mathrm{spl}}))
            \gets\textsc{RGibbsScan}(\mathbf{L}_R,\dots)$ in sampling mode.
    \State Form $\mathcal{P}^{\mathrm{spl}}$ from $\mathbf{L}^{\mathrm{fin}}$;
    set $\log q_{\mathrm{prop}}(\mathcal{P}^{\mathrm{spl}}\to\mathcal{P})=0$.
    \State Compute $\log p_{\EPA}(\mathcal{P})$ and $\log p_{\EPA}(\mathcal{P}^{\mathrm{spl}})$.
    \State $\log R_{\mathrm{spl}}\gets$ \eqref{eq:logR-split}.
    \State Draw $u\sim\Unif(0,1)$; if $\log u<\log R_{\mathrm{spl}}$, accept:
    $\mathcal{P}\gets\mathcal{P}^{\mathrm{spl}}$.
\Else
    \State Form $\mathcal{P}^{\mathrm{mrg}}$;
    set $\log q_{\mathrm{prop}}(\mathcal{P}\to\mathcal{P}^{\mathrm{mrg}})=0$.
    \State Initialize $\mathbf{L}_0^{\mathrm{rev}}\gets\{S_u,S_v\}$.
    \For{$r=1,\dots,R$}
        \State $(\mathbf{L}_r^{\mathrm{rev}},\cdot)\gets\textsc{RGibbsScan}(\mathbf{L}_{r-1}^{\mathrm{rev}},\dots)$ in sampling mode.
    \EndFor
    \State Run \textsc{RGibbsScan} from $\mathbf{L}_R^{\mathrm{rev}}$ in
    \emph{scoring mode} forcing $L_{q^*}^{\mathrm{target}}=L_{q^*}^{\mathcal{P}}$,
    accumulating $\log q_{\mathrm{prop}}(\mathcal{P}^{\mathrm{mrg}}\to\mathcal{P})$.
    \State Compute $\log p_{\EPA}(\mathcal{P})$ and $\log p_{\EPA}(\mathcal{P}^{\mathrm{mrg}})$.
    \State $\log R_{\mathrm{mrg}}\gets$ \eqref{eq:logR-merge}.
    \State Draw $u\sim\Unif(0,1)$; if $\log u<\log R_{\mathrm{mrg}}$, accept:
    $\mathcal{P}\gets\mathcal{P}^{\mathrm{mrg}}$.
\EndIf
\State \Return $\mathcal{P}$.
\end{algorithmic}
\end{algorithm}

\begin{algorithm}[H]
\caption{One full MCMC sweep (EPA--LSIRM).}
\small
\begin{algorithmic}[1]
\State \textbf{Input:} data $Y$, hyperparameters, proposal scales,
$M_\sigma$, $M_{\mathrm{SM}}$, $R$.
\State Update layer-specific nuisance $\xi_\ell$ (e.g.\ $\gamma_\ell$,
$\sigma_0^2$, $\lambda_{ij}^{(\ell)}$, $\{\kappa_j\}$) by their existing
updates.
\For{$i=1,\dots,n$}
    \State Propose $a_i'$ and accept/reject via \eqref{eq:a-mh-ratio}.
\EndFor
\For{$\ell=1,\dots,L$}
    \For{$i=1,\dots,n$}
        \State Propose $\alpha_i^{(\ell)\prime}$ and accept/reject via
        \eqref{eq:alpha-mh-ratio}.
    \EndFor
    \For{$j=1,\dots,P_\ell$}
        \State Propose $\beta_j^{(\ell)\prime}$ via \eqref{eq:beta-mh-ratio}
        (or thresholds for ordinal layers, with ordering check).
        \State Propose $b_j^{(\ell)\prime}$ and accept/reject via
        \eqref{eq:b-mh-ratio}, including the EPA prior contribution.
    \EndFor
    \State Update $\sigma_{\alpha,\ell}^2$ via \eqref{eq:sigma-alpha-update}.
\EndFor
\State Form global item positions $z_q=b_{j(q)}^{(\ell(q))}$ and refresh
the cached pairwise similarity matrix $\Lambda$.
\For{$q=1,\dots,P$ in random order}
    \State Update $\mathcal{P}$ at $q$ by single-item EPA Gibbs
    \eqref{eq:single-item-gibbs}.
\EndFor
\For{$m=1,\dots,M_{\mathrm{SM}}$}
    \State Run one Jain--Neal split--merge step (Algorithm~2).
\EndFor
\For{$m=1,\dots,M_\sigma$}
    \State Random-swap MH update of $\sigma$ via \eqref{eq:sigma-mh-ratio}.
\EndFor
\State Log-scale MH update of $\alpha$ via \eqref{eq:alpha-hyp-mh-ratio}.
\State Log-scale MH update of $\tau$ (analogous).
\If{iteration is after burn-in}
    \State Update posterior similarity counts
    $C_{qr}^{\mathrm{raw}}\leftarrow C_{qr}^{\mathrm{raw}}+\ind(q\sim_{\mathcal{P}}r)$.
    \State Store $\mathcal{P}$, $K=|\mathcal{P}|$, $z$, and selected LSIRM
    parameters.
\EndIf
\end{algorithmic}
\end{algorithm}

\section{Differences from v11 at a glance}

\begin{center}
\begin{tabular}{p{0.27\textwidth} p{0.34\textwidth} p{0.33\textwidth}}
\toprule
Component & v11 (latent-position GMM) & v12 (EPA partition) \\
\midrule
Item position prior &
$b_j^{(\ell)}\mid c_q\sim N_d(\mu_{c_q},\Sigma_{c_q})$ &
$b_j^{(\ell)}\sim N_d(0,\sigma_b^2 I_d)$, free \\
Clustering target &
labels $c_q\in\{1,\dots,K^\star\}$, $\rho\sim\Dir$, $(\mu_l,\Sigma_l)\sim$NIW &
partition $\mathcal{P}$ with EPA pmf, no component parameters \\
Cluster scale prior &
NIW with Wishart hyperprior on $S_0$ &
none (no NIW); EPA uses sampled $\tau$ on raw squared distances \\
Auxiliary state &
$c$, $\rho$, $\{\mu_l,\Sigma_l\}$, $S_0$ &
$\mathcal{P}$, $\sigma\in\Perm(P)$, $(\alpha,\tau)$ \\
Allocation update &
NIW collapsed predictive Gibbs (Student-$t$ weights) &
EPA single-item Gibbs (\ref{eq:single-item-gibbs}) \\
Split--merge target &
$\pi(c\mid z)\propto\prod_l\Gamma(n_l+e_0)m(z_{\mathcal{J}_l})$ &
$p_{\EPA}(\mathcal{P}\mid\sigma,z,\alpha,\delta,\tau)$ \\
Split--merge variant &
Jain--Neal conjugate (NIW marginal cancels into $m$) &
Jain--Neal nonconjugate; LSIRM likelihood cancels by \eqref{eq:joint-decouple} \\
$b$ MH ratio includes &
LSIRM + mixture-prior term &
LSIRM + free Gaussian + EPA-prior term (\ref{eq:b-mh-ratio}) \\
Variance-inflation hack on $b$ &
optional (v11 §\ldots) &
not needed: no cluster covariance pulls $b$ \\
\bottomrule
\end{tabular}
\end{center}

\end{document}
